\begin{enumerate}
\def\labelenumi{\arabic{enumi}.}
\item
  \begin{enumerate}
  \def\labelenumii{\arabic{enumii}.}
  \item
    A continuación hablaremos sobre como son en general las álgebras de
    Boole, funda\-mentalmente las finitas, y veremos que efectivamente
    podemos llegar a teoremas que nos dicen de forma muy concreta cuales
    son estas álgebras de Boole. Vamos a ver for\-mas de generarlas y
    cuestiones parecidas.

    \begin{enumerate}
    \def\labelenumiii{\arabic{enumiii}.}
    \item
      En el caso que el cardinal de \[B\]sea finito,
      \[2 \mid {({\# B})}\]. Para demostrarlo solo hay que darse cuenta
      que \[{B = \cup_{x \in B}}{\{{x,\overline{x}}\}}\], que
      \[{({{x \neq {y \land \overline{x}}} \neq y})}\Rightarrow{({{{\{{x,\overline{x}}\}} \cap {\{{y,\overline{y}}\}}} = \varnothing})}\]
      y que \[\forall{x \in B}\#{{\{{x,\overline{x}}\}} = 2}\] y así
      cuando\[B\]sea finito, su cardinal será un múlti\-plo de 2.
    \item
      Definición:

      \[\forall x,{y \in B}{x \leq y}\Leftrightarrow{{x \cdot y} = x}\Leftrightarrow{{x + y} = y}\]
    \item
      \[\left\langle {B, \leq} \right\rangle\text{es un}\mathit{orden}\].

      \begin{enumerate}
      \def\labelenumiv{\arabic{enumiv}.}
      \item
        Reflexiva: \[\forall{x \in B}{x \leq x}\]

        \begin{enumerate}
        \def\labelenumv{\arabic{enumv}.}
        \tightlist
        \item
          \[\forall{x \in B}{{x \cdot x} = x}\]
        \item
          \[\forall{x \in B}{x \leq x}\]
        \end{enumerate}
      \item
        Antisimétrica:
        \[\forall x,{y \in B}{{x \leq {y \land y}} \leq x}\Rightarrow{x = y}\]

        \begin{enumerate}
        \def\labelenumv{\arabic{enumv}.}
        \tightlist
        \item
          \[x,{y \in B}{{x \leq {y \land y}} \leq x}\]
        \item
          \[x,{y \in B}{{{x \cdot y} = {{y \land x} \cdot y}} = x}\]
        \item
          \[x,{y \in B}{x = y}\]
        \end{enumerate}
      \item
        Transitiva:
        \[\forall x,y,{z \in B}{{x \leq {y \land y}} \leq z}\Rightarrow{x \leq z}\]

        \begin{enumerate}
        \def\labelenumv{\arabic{enumv}.}
        \tightlist
        \item
          \[{x \leq {y \land y}} \leq z\]
        \item
          \[{{x \cdot y} = {{x \land y} \cdot z}} = y\]
        \item
          \[{{x \cdot y} \cdot z} = {x \cdot z}\]
        \item
          \[{x \cdot y} = {x \cdot z}\]
        \item
          \[x = {x \cdot z}\]
        \item
          \[x \leq z\]
        \end{enumerate}
      \end{enumerate}
    \item
      Definición:\[\forall x,{y \in B}{x \geq y}\Leftrightarrow{{x \cdot y} = y}\Leftrightarrow{{x + y} = x}\]
    \item
      \[\left\langle {B, \geq} \right\rangle\text{es un}\mathit{orden}\].

      \begin{enumerate}
      \def\labelenumiv{\arabic{enumiv}.}
      \item
        Reflexiva: \[\forall{x \in B}{x \geq x}\]

        \begin{enumerate}
        \def\labelenumv{\arabic{enumv}.}
        \tightlist
        \item
          \[\forall{x \in B}{{x + x} = x}\]
        \item
          \[\forall{x \in B}{x \geq x}\]
        \end{enumerate}
      \item
        Antisimétrica:
        \[\forall x,{y \in B}{{x \geq {y \land y}} \geq x}\Rightarrow{x = y}\]

        \begin{enumerate}
        \def\labelenumv{\arabic{enumv}.}
        \tightlist
        \item
          \[x,{y \in B}{{x \geq {y \land y}} \geq x}\]
        \item
          \[x,{y \in B}{{{x + y} = {{y \land x} + y}} = x}\]
        \item
          \[x,{y \in B}{x = y}\]
        \end{enumerate}
      \item
        Transitiva:
        \[\forall x,y,{z \in B}{{x \geq {y \land y}} \geq z}\Rightarrow{x \geq z}\]

        \begin{enumerate}
        \def\labelenumv{\arabic{enumv}.}
        \tightlist
        \item
          \[{x \geq {y \land y}} \geq z\]
        \item
          \[{{x + y} = {{x \land y} + z}} = y\]
        \item
          \[{{x + y} + z} = {x + z}\]
        \item
          \[{x + y} = {x + z}\]
        \item
          \[x = {x + z}\]
        \item
          \[x \geq z\]
        \end{enumerate}
      \end{enumerate}
    \item
      Definición:\[\mathit{atom}{(x)}\overset{\text{def}}{\Leftrightarrow}{\left\lbrack {x \in (1,0)_{B}} \right\rbrack \land \left\lbrack {\forall{y \in (1,0)_{B}}\left( {\left( {{x \cdot y} = x} \right) \vee \left( {{x \cdot y} = 0} \right)} \right)} \right\rbrack}\]
    \item
      \[\forall x,{y \in B}\mathit{atom}{{(x)} \land \mathit{atom}}{(y)}\Rightarrow{{x \cdot y} = 0}\].
    \item
      Definición:\[\mathit{hatom}{(x)}\overset{\text{def}}{\Leftrightarrow}{\left\lbrack {x \in \left( {B \smallsetminus {\{ 0,1\}}} \right)} \right\rbrack \land \left\lbrack {\forall{y \in \left( {B \smallsetminus {\{ 0,1\}}} \right)}\left( {\left( {{x + y} = x} \right) \vee \left( {{x + y} = 1} \right)} \right)} \right\rbrack}\]
    \item
      \[\forall x,{y \in B}\mathit{hatom}{{(x)} \land \mathit{hatom}}{(y)}\Rightarrow{{x + y} = 1}\]
    \item
      Definición\[{x \in B}\Rightarrow{\lbrack{x,0}\rbrack}_{B}{: = {\{{{y \in B} \mid {y \leq x}}\}}}\]
    \item
      \[\#{{\lbrack{x,0}\rbrack}_{B} = 1}\Leftrightarrow{{\lbrack{x,0}\rbrack}_{B} = {\{ 0\}}}\Leftrightarrow{x = 0}\]
    \item
      Definición\[{x \in B}\Rightarrow{\lbrack{1,x}\rbrack}_{B}{: = {\{{{y \in B} \mid {y \geq x}}\}}}\]
    \item
      \[\#{{\lbrack{1,x}\rbrack}_{B} = 1}\Leftrightarrow{{\lbrack{1,x}\rbrack}_{B} = {\{ 1\}}}\Leftrightarrow{x = 1}\]
    \item
      Definición\[x,{y \in B}{x \geq y}\Rightarrow\left\lbrack {x,y} \right\rbrack_{B}{: = {\{{{z \in B} \mid {{y \leq {z \land x}} \geq z}}\}}}\]
    \item
      Definición\[x,{y \in B}{x \geq y}\Rightarrow\left\lbrack {x,y} \right)_{B}{: = {\{{{z \in B} \mid {{{y \leq {z \land x}} \geq {z \land z}} \neq y}}\}}}\]
    \item
      Definición\[x,{y \in B}{x \geq y}\Rightarrow\left( {x,y} \right\rbrack_{B}{: = {\{{{z \in B} \mid {{{y \leq {z \land x}} \geq {z \land z}} \neq x}}\}}}\]
    \item
      Definición\[x,{y \in B}{x \geq y}\Rightarrow\left( {x,y} \right)_{B}{: = {\{{{z \in B} \mid {{{{y \leq {z \land x}} \geq {z \land z}} \neq {x \land z}} \neq y}}\}}}\]
    \item
      \[{x \in B}\Rightarrow{{{\lbrack{x,0}\rbrack}_{B} \cap {\lbrack{1,x}\rbrack}_{B}} = {\{ x\}}}\]
    \item
      \[{x \in B}\Rightarrow{{{\lbrack{x,0}\rbrack}_{B} \cap {\lbrack{\overline{x},0}\rbrack}_{B}} = {\{ 0\}}}\]
    \item
      \[{x \in B}\Rightarrow{{{\lbrack{1,x}\rbrack}_{B} \cap {\lbrack{1,\overline{x}}\rbrack}_{B}} = {\{ 1\}}}\]
    \item
      \[\forall{x \in B}\forall{y \in \left( {x,0} \right)_{B}}{{({{\lbrack{y,0}\rbrack}_{B} \subset {\lbrack{x,0}\rbrack}_{B}})} \land {({{\lbrack{y,0}\rbrack}_{B} \neq {\lbrack{x,0}\rbrack}_{B}})}}\]
    \item
      \[\forall{x \in B}\forall{y \in \left( {1,x} \right)_{B}}{{({{\lbrack{1,y}\rbrack}_{B} \subset {\lbrack{1,x}\rbrack}_{B}})} \land {({{\lbrack{1,y}\rbrack}_{B} \neq {\lbrack{1,x}\rbrack}_{B}})}}\]
    \item
      Definición\[\mathit{Atom}B{: = {\{{{x \in (1,0)_{B}} \mid \mathit{atom}{(x)}}\}}}\]
    \item
      Definición\[\mathit{Hatom}B{: = {\{{{x \in (1,0)_{B}} \mid \mathit{hatom}{(x)}}\}}}\]
    \item
      \[\mathit{Atom}{{(B_{2})} = \mathit{Hatom}}{{(B_{2})} = \varnothing}\]
    \item
      Definición
      \[\left\lbrack B \right){: = {\{{{A \subset B} \mid \exists{x \in B}{A = \left\lbrack {x,0} \right)_{B}}}\}}}\]
    \item
      \[\left\lbrack B \right) \neq \varnothing\]. Pues es un álgebra de
      cardinal mayor o igual que 2 y existe al menos \[\{ 1\}\].
    \item
      \[{\langle{\left\lbrack B \right), \supseteq}\rangle}\mathit{es}\mathit{un}\mathit{orden}\].
    \item
      \[B\mathit{finito}\forall{P \subset \left\lbrack B \right)}{\langle{P, \supseteq}\rangle}\mathit{orden}\mathit{total}\Rightarrow\exists!{x \in \underset{X \in P}{\cap}}X\mathit{atom}(x)\]

      \begin{enumerate}
      \def\labelenumiv{\arabic{enumiv}.}
      \tightlist
      \item
        Prueba:
      \item
        \[{\langle{P, \supseteq}\rangle}\mathit{es}\mathit{un}\mathit{orden}\mathit{total}\]
      \item
        \[\forall X,{Y \in P}{X \neq Y}\Rightarrow{{X \supset {Y \vee Y}} \supset X}\]
      \item
        \[\forall X,{Y \in P}{X \neq Y}\Rightarrow{\left( {\left( {Y \supset X} \right) \vee \left( {Y \supset X} \right)} \right) \land \left( {\left( {{Y \cap X} = X} \right) \vee \left( {{Y \cap X} = Y} \right)} \right)}\]
      \item
        \[\forall X,{Y \in P}{X \neq Y}\Rightarrow{\left( {\left( {{Y \cap X} = X} \right) \vee \left( {{Y \cap X} = Y} \right)} \right) \land \left( {{X \neq {{\{ 0\}} \land Y}} \neq {\{ 0\}}} \right)}\]
      \item
        \[\forall X,{Y \in P}\left( {X \neq Y} \right)\Rightarrow\left( {{Y \cap X} \neq {\{ 0\}}} \right)\]
      \item
        \[\underset{X \in P}{\cap}{X \neq {{\{ 0\}} \land \underset{X \in P}{\cap}}}{X \supset {\{ 0\}}}\]
      \item
        \[\underset{Y \in P}{\cap}Y \supsetneq {\{ 0\}}\]
      \item
        \[\underset{Y \in P}{\cap}{Y \supseteq {\{{0,x}\}}}\]
      \item
        \[\underset{Y \in P}{\cap}{Y \supseteq {\lbrack{x,0}\rbrack}_{B}}\]
      \item
        \[\underset{Y \in P}{\cap}{{Y \supseteq {\lbrack{x,0}\rbrack}_{B}} \supseteq {\{{0,x}\}}}\]
      \item
        \[{\lbrack{x,0}\rbrack}_{B} \in P\]
      \item
        \[\exists{X \in P}\exists{x \in X}{\underset{Y \in P}{\cap}{Y = {\lbrack{x,0}\rbrack}_{B}}}\]
      \item
        \[\exists!{x \in B}{\underset{Y \in P}{\cap}{Y = {\lbrack{x,0}\rbrack}_{B}}}\]
      \item
        \[\exists{x \in {\underset{Y \in P}{\cap}Y{{\lbrack{x,0}\rbrack}_{B} = {\{{0,x}\}}}}}\]
      \item
        \[\exists{x \in {\underset{Y \in P}{\cap}Y\forall{X \in P}{{X \supseteq {\lbrack{x,0}\rbrack}_{B}} = {\{{0,x}\}}}}}\]
      \item
        \[\forall{P \subset \left\lbrack B \right)}{\langle{P, \supseteq}\rangle}\mathit{orden}\mathit{total}\Rightarrow\exists!{x \in \underset{X \in P}{\cap}}X\mathit{atom}(x)\]
      \end{enumerate}
    \item
      \[B\mathit{finito}\Rightarrow\mathit{Atom}{B \neq \varnothing}\].
      Desde 57 es inmediato.
    \item
      Ahora vamos a construir una función inyectiva del álgebra de Boole
      de las partes de los átomos de B (si este es finito) en el álgebra
      de Boole B. Así cuando menos sa\-bremos que podemos interpretar
      este álgebra de las partes de un conjunto de los áto\-mos de B
      como un subálgebra de la que estamos estudiando.

      La función \[\varphi\]que vamos a definir va a quedar
      completamente definida en la fór\-mula que sigue. Tendremos que
      mostrar que está bien definida, que es inyectiva, que
      \[\varphi{{({x \cup y})} = \varphi}{{(x)} + \varphi}{(y)}\], esto
      es que respeta la suma booleana en \[B\]que viene como unión de
      conjuntos desde \[\wp\left( {\mathit{Atom}{(B)}} \right)\], que
      \[\varphi{{({x \cap y})} = \varphi}{{(x)} \cdot \varphi}{(y)}\],
      esto es que respeta el producto booleano en \[B\]que viene como
      intersección de conjun\-tos desde
      \[\wp\left( {\mathit{Atom}{(B)}} \right)\] , y aunque ya no sería
      necesario, también veremos que
      \[\varphi{\left( {\mathit{Atom}{{(B)} \smallsetminus x}} \right) = \overline{\varphi(x)}}\]
      . Así quedará clara la relación entre ambas álgebras.

      \begin{enumerate}
      \def\labelenumiv{\arabic{enumiv}.}
      \tightlist
      \item
        \[\begin{matrix}
        {n{: = \#}\left( {\mathit{Atom}B} \right)} \\
        {{n \leq m}{: = \#}\left( B \right)} \\
        {\left\lbrack {1,n} \right\rbrack_{\mathbb{N}}{: = {\{{1,2,\ldots,n}\}}}} \\
        {\varphi:\wp{\left( {\mathit{Atom}\left( B \right)} \right)\rightarrow B}} \\
        {{\varphi{(x)}}{: =}\begin{Bmatrix}
        0 & \Leftarrow & {{x = \varnothing} = {\{\}}} \\
        {x_{1}'} & \Leftarrow & {x = {\{{x_{1}'}\}}} \\
        {{\sum\limits_{\substack{k \in I \\ I \subset {\lbrack{1,n}\rbrack}_{\mathbb{N}}}}x_{k}}'} & \Leftarrow & {x = {\{{{x_{i}'} \mid {{i \in I} \subset \left\lbrack {1,n} \right\rbrack_{\mathbb{N}}}}\}}} \\
        {x_{1}{{' + \ldots} + x_{i - 1}}{' + x_{i}}{{' + \ldots} + x_{n}}'} & \Leftarrow & {{x = \mathit{Atom}}{{(B)} \smallsetminus {\{{x_{i}'}\}}}} \\
        1 & \Leftarrow & {{x = \mathit{Atom}}{(B)}}
        \end{Bmatrix}}
        \end{matrix}\].
      \end{enumerate}
    \item
      \[\#{B_{a} = \#}{B_{b} = 2}\Rightarrow{B_{a} \simeq B_{b}}\]. Con
      la misma \[\varphi\] anterior.
    \item
      \[\#{B_{a} = \#}{B_{b} = 4}\Rightarrow{B_{a} \simeq B_{b}}\]. Con
      la misma \[\varphi\] anterior.
    \item
      \[\#{B > 4}\Rightarrow\mathit{Atom}{B \cap \mathit{Hatom}}{B = \varnothing}\].
      Lo mejor sería demostrar que cualquier cadena completa saturada de
      1 a 0 tiene una longitud (número de elementos) siempre igual al
      \[\#\mathit{Atom}{(B)}\]. De ahí se sigue que si el cardinal es el
      dicho, tendríamos más de 3 niveles, diferenciándose siempre los
      átomos y los hiperátomos.
    \item
      Para cada uno de los cardinales de \[B\], cuando son finitos,
      existe una estructura no solo de anillo conmutativo con unidad
      como en 34, sino también de cuerpo. La po\-demos encontrar
      explícitamente en el álgebra de las partes de un conjunto finito.
      Sólo nos queda ver que \[\varphi\]es sobreyectivo. Tenemos que
      \[\varphi{\left( {\mathit{Atom}{(B)}} \right) = B}\]. Así
      \[\varphi\]pasa a ser un isomorfismo de álgebras de Boole: esto
      es, en lo que a la estructu\-ra de álgebra de Boole se refiere,
      haciendo abstracción de las operaciones \[{} + {}\] y
      \[{} \cdot {}\]concretas y los elementos concretos,
      \[{\langle{B,0,1, + , \cdot}\rangle} \simeq {\langle{\wp\left( {\mathit{Atom}{(B)}} \right),\varnothing,\mathit{Atom}\left( B \right), \cup , \cap}\rangle}\].
    \item
      De 74 se deduce que si \[B\]es un conjunto finito,
      \[\exists{n \in \mathbb{N}}\#{B = 2^{\mathbf{\mathrm{n}}}}\].
    \end{enumerate}
  \end{enumerate}
\end{enumerate}
